\section{State of the Art}
\label{StateOfTheArt}

Παραδοσιακά, το πρόβλημα της αναγνώρισης της ανθρώπινης ομιλίας αξιοποιώντας οπτικά στοιχεία (βίντεο της περιοχής των χειλιών) αντιμετωπιζόταν ως μία δυσκολότερη εκδοχή της ακουστικής ASR (Automatic Speech Recognition). Χωρίς ηχητικά στοιχεία, το γλωσσικό περιεχόμενο πρέπει να αντληθεί χρησιμοποιώντας αποκλειστικά τις κινήσεις των χειλιών και των σαγονιών, με την πληροφορία που αφορά κυρίως τον τρόπο εκφοράς του κάθε φωνήματος να χάνεται. Οι περιορισμοί αυτοί οδήγησαν, αρχικά, σε μελέτες που εστίαζαν στην ταξινόμηση συγκεκριμένων φωνημάτων και visemes από στατικές εικόνες ή μικρά βίντεο σε αντίθεση με τη χρήση ολόκληρων λέξεων ή φράσεων. Προφανώς η αντιμετώπιση αυτή υπέφερε από διάφορους περιορισμούς, κυρίως του μικρού μεγέθους του dataset που χρησιμοποιούταν, το οποίο ήταν ανίκανο να περιγράψει πλήρως τη φύση της αλληλουχίας των φωνημάτων στην ανθρώπινη ομιλία. Μερικά παραδείγματα τέτοιων μελετών είναι η \cite{Yuhas1989IntegrationOA}, στην οποία συνδυάζονται οπτικοακουστικά στοιχεία για να βελτιωθεί η αναγνώριση ομιλίας με ακουστικό θόρυβο, η \cite{10.1121/1.401235}, στην οποία συγκρίνεται η επίδοση νευρωνικού δικτύου και ανθρώπων στην αναγνώριση φωνηέντων από στατικές εικόνες και η \cite{inbook}, στην οποία γίνεται ταξινόμηση σε visemes.

Με την έλευση του deep learning, άρχισαν να αναπτύσσονται end-to-end αρχιτεκτονικές αξιοποιώντας τα 3D CNNs και τα RNNs. Το LipNet \cite{assael2016lipnetendtoendsentencelevellipreading} αξιοποίησε τη συνέλιξη στις 3 διαστάσεις με bi-directional GRUs και τη CTC loss (βλ. \ref{MachineLearningBackground}) ώστε να ταξινομήσει βίντεο της περιοχής του στόματος απ'ευθείας σε αλληλουχίες χαρακτήρων, πετυχαίνοντας υψηλή ακρίβεια στο περιορισμένο GRID dataset, το οποίο χρησιμοποιείται και στα πλαίσια αυτής της διπλωματικής. Λίγο αργότερα, οι Chung et. al δημιούργησαν το Watch, Listen, Attend and Spell \cite{Chung17}, ένα μοντέλο sequence-to-sequence με ξεχωριστό encoder για ``Watch'' (οπτικά στοιχεία, 3D-CNN + LSTM) και ``Listen'' (ακουστικά στοιχεία, MFCC + LSTM), οι οποίοι τροφοδοτούσαν έναν LSTM decoder βασισμένο στον μηχανισμό του attention. Με τη δουλειά αυτή καθιερώθηκαν τα sentence-level, in-the-wild datasets που χρησιμοποιούνται ακόμη και σήμερα ως benchmarks - τα LRW, LRS2 και LRS3, \cite{Afouras18c}, \cite{Chung16}, \cite{Chung17}, \cite{Chung17a}. Οι δουλειές αυτές απέδειξαν, επίσης, πόσο πιο δύσκολο ήταν πραγματικά το πρόβλημα σε μη περιορισμένο περιβάλλον: ακόμα και μοντέλα με μεγάλη ακρίβεια σε επίπεδο λέξεων στο GRID υπέφεραν σε αυτά τα in the wild datasets.

Μετέπειτα εργασίες εστίασαν στη βελτίωση του temporal encoder και όχι της αντικειμενικής συνάρτησης ή του ορισμού του προβλήματος. Οι Afouras et. al \cite{Afouras18c} σύγκριναν back-ends βασισμένα σε συνελικτικά νευρωνικά δίκτυα, BiLSTM ή Transformer στο πρόβλημα του lip reading, καταλήγοντας ότι το self-attention πετύχαινε τελικά τις βέλτιστες τιμές ακρίβειας σε επίπεδο λέξης. Έπειτα, σύγκριναν τις επιδόσεις ενός Transformer decoder με CTC και με sequence-to-sequence βάση, τόσο για είσοδο αποκλειστικά οπτική όσο και για οπτικοακουστική. Η περίοδος αυτή καθιέρωσε δύο αρχιτεκτονικές επιλογές οι οποίες έμειναν σταθερές (ως επί το πλείστον) στις περισσότερες μετέπειτα εργασίες, όπως και σε αυτήν: ένα front-end που αξιοποιεί τόσο τη χρονική όσο και τις χωρικές διαστάσεις του βίντεο (3D CNN ή Transformer), διατηρώντας πληροφορία σε επίπεδο frame, και έναν decoder βασιμένο είτε στο CTC είτε στον μηχανισμό του attention.

Πρόσφατες εξελίξεις έχουν αλλάξει τον τρόπο με τον οποίο αξιοποιούνται οι ροές δεδομένων των οπτικών αλλά και των ακουστικών στοιχείων. Το AV-HuBERT \cite{shi2022learningaudiovisualspeechrepresentation} συνέχισε τη δουλειά του HuBERT \cite{hsu2021hubertselfsupervisedspeechrepresentation} με το masked-prediction objective, μεταφέροντάς το από το ακουστικό στο οπτικοακουστικό πρόβλημα αναγνώρισης ομιλίας. Το προτεινόμενο μοντέλο της εργασίας αυτής μαθαίνει μία κοινή αναπαράσταση των οπτικοακουστικών στοιχείων, χρησιμοποιώντας multimodal clustering και modality dropout, πετυχαίνοντας τελικά βελτιωμένη επίδοση σε σχέση με το τότε state-of-the-art, ακόμα κι αν χρησιμοποιούσε πολύ λιγότερα labelled δεδομένα. Αντιθέτως, το Auto-AVSR \cite{Ma_2023} επέκτεινε το supervision χρησιμοποιώντας ήδη υπάρχοντα ASR μοντέλα ώστε να παράξει αυτόματα labels για πολλές ώρες οπτικοακουστικού υλικού. Άλλες εργασίες, όπως το Lip2Vec \cite{djilali2023lip2vecefficientrobustvisual}, αποτυπώνουν τα features εξαγόμενα από οπτικά στοιχεία σε προεκπαιδευμένους latent χώρους από ακουστικά μοντέλα. Σε κάθε περίπτωση, η συνένωση των οπτικών και των ακουστικών στοιχείων μέσω cross-modal ή cross-attention μηχανισμών αποτελεί πρακτική που εφαρμόζεται σε μεγάλο βαθμό, όπως και στη συγκεκριμένη εργασία, στην περίπτωση του audiovisual μοντέλου που παρουσιάζεται στην ενότητα \ref{AudiovisualModel}.

Για να αντιμετωπιστεί η ασάφεια που προκύπτει στο πρόβλημα του lip reading εξαιτίας των visemes, πρόσφατες μελέτες μεταβάλλουν το τι ακριβώς καλείται να προβλέψει ο visual encoder και όχι πώς να το προβλέψει. Συγκεκριμένα, αντί να αποτυπώνεται ένα βίντεο απ'ευθείας σε αλληλουχία χαρακτήρων ή λέξεων, πολλά συστήματα αξιοποιούν ένα ενδιάμεσο στάδιο πρόβλεψης φωνημάτων, με τη λογική ότι η αλληλουχία φωνημάτων είναι τόσο ένας πιο εύκολος στόχος ως αντικείμενο πρόβλεψης από οπτικά στοιχεία, όσο και ένα ικανοποιητικό ενδιάμεσο στάδιο στο οποίο μπορεί να εφαρμοστεί περαιτέρω γλωσσικό context, ώστε να επιλυθεί η ασάφεια αυτή. Αυτή τη λογική ακολουθεί το VALLR \cite{thomas2026vallrvisualasrlanguage}, το οποίο αποτελεί πηγή έμπνευσης για μεγάλο τμήμα της συγκεκριμένης εργασίας. Το VALLR χρησιμοποιεί έναν Video Transformer εκπαιδευμένο με CTC ώστε να προβλέπει αλληλουχίες φωνημάτων από το βίντεο της εισόδου, οι οποίες σε ξεχωριστό στάδιο εισέρχονται σε ένα κατάλληλα fine-tuned LLM που παράγει τις προβλεπόμενες λέξεις και τελικά φράσεις, αξιοποιώντας γλωσσικό παρακείμενο σε επίπεδο φράσης στο οποίο δεν μπορεί να έχει πρόσβαση ο visual encoder. Η εκπαίδευση γίνεται σε δύο βήματα, χωρίζοντας το οπτικά αμφίσημο πρόβλημα της αναγνώρισης των φωνημάτων από το πρόβλημα αποκωδικοποίησης των εν λόγω φωνημάτων σε λέξεις (κάτι που με γνώση της γλώσσας που διαθέτει ένα LLM, γίνεται με μεγάλη επιτυχία). Το προτεινόμενο μοντέλο πετυχαίνει state of the art επιδόσεις στα προαναφερθέντα in the wild datasets, χρησιμοποιώντας λιγότερα labelled δεδομένα εκπαίδευσης από άλλα πλήρως supervised μοντέλα ή μοντέλα με self-supervised pre-training, ενώ δεν απαιτεί οπτικοακουστικό pre-training μεγάλης κλίμακας όπως το AV-HuBERT.

Το κεντρικό νόημα της μελέτης του VALLR, ότι δηλαδή η επίβλεψη σε επίπεδο φωνήματος, σε αντίθεση με την πρόβλεψη χαρακτήρων ή λέξεων, επιτρέπει σε μικρούς visual encoders να επιτύχουν υψηλή ακρίβεια στις προβλέψεις τους, αποτελεί την αρχή για τη συγκεκριμένη εργασία. Η εξαγωγή των φωνημάτων σε επίπεδο frame, η αντικειμενική συνάρτηση που περιλαμβάνει γνώση των visemes, η συνένωση των οπτικών και των ακουστικών features με βάση την αξιοπιστία τους, που αναλύονται στα επόμενα κεφάλαια, λειτουργούν ως περαιτέρω εξερεύνηση της αφοσίωσης στα φωνήματα για το πρόβλημα της αναγνώρισης ομιλίας. Έτσι, μελετάται το πώς εξάγονται και ευθυγραμμίζονται τα φωνήματα σε επίπεδο frame, πώς πρέπει να αντιμετωπιστούν οι κλάσεις φωνημάτων που ανήκουν σε visemes και άρα παρουσιάζουν υψηλή σύγχυση και πώς μπορεί να αξιοποιηθεί ένα ακουστικό σήμα (ακόμα και θορυβώδες) σε συνδυασμό με το οπτικό σήμα για τη βελτίωση των προβλέψεων.